% introduction.tex — §1

\section{Introduction}
\label{sec:intro}

The Dark Energy Spectroscopic Instrument (DESI) has measured baryon acoustic oscillations (BAO) across $0 < z < 2.3$ with unprecedented precision \citep{DESI_DR2_BAO}. Combined with cosmic microwave background (CMB) data and Type~Ia supernovae (SNe~Ia), the DESI collaboration reports a $2.8$--$3.9\sigma$ preference (depending on SN dataset) for evolving dark energy in the \wowa parameterization \citep{DESI_DR2_cosmo}. What drives this preference? We show that it arises primarily from what could be interpreted as a $\sim 2\sigma$ discrepancy in \omh between the CMB and low-redshift distance data, rather than from any dark energy signal.

\claim{w0wa_is_c0}{The DESI \wowa preference comes entirely from \czero (\omh), not from the dark energy direction.}

There are many ways to search for deviations from \LCDM. Here we take \LCDM as a well-established baseline whose parameters are constrained by the CMB (we use ACT DR6; \citealt{ACT_DR6}) to a small allowed range. The question we ask is: \emph{how large is the space of deviations that low-redshift distance probes can actually detect?}

The answer is simpler than one might expect. The CMB-allowed \LCDM parameter range maps into a low-dimensional subspace of distance predictions. For both BAO and SN measurements, the effective dimensionality of this map is one: a single SVD direction $V_0$ along which the CMB allows essentially all of the model motion. We write \czero for the projection of the data onto $V_0$. The matrix being decomposed has one row per CMB chain sample and one column per whitened observable (Section~\ref{sec:svd_method}); its leading singular value exceeds the next by a factor of $\sim 25$ for BAO. The result depends only on the CMB parameter uncertainties and the low-redshift measurement covariance --- not on the low-redshift data values.

We can then ask the same question for extensions beyond \LCDM: does a given extension (dynamical dark energy, spatial curvature, etc.) open a genuinely new measurable direction in observable space, beyond $V_0$? Again, this is answerable from the theory and covariance alone.

The result is a set of predetermined directions onto which we project the data. Searching for deviations reduces to reading off a few numbers (tensions in units of $\sigma$, one per mode).

\textbf{Preview of results.} All tension concentrates in a single universal direction $V_0$: BAO shows $+2.2\sigma$ in \czero against the CMB prediction, while the three SN compilations are individually mild (DES-Dovekie $-0.8\sigma$, Pantheon+ $-1.1\sigma$, Union3 $-1.7\sigma$) and do not have enough constraining power in \czero to confirm or refute the BAO signal. The independent dark-energy direction $V_1$, which becomes measurable in the \wowa extension, shows only $+1.2\sigma$ for BAO --- consistent with a cosmological constant. The only other extension that opens a measurable new direction is spatial curvature, and only for BAO.

Among the other extensions we considered, the phenomenological lensing rescaling $A_\mathrm{lens}$ --- a multiplicative factor on the smoothing of CMB acoustic peaks by gravitational lensing, with $A_\mathrm{lens} = 1$ in \LCDM \citep{Calabrese2008} --- is notable: it reduces the tension by shifting the predicted \czero mean and increasing its variance, reminding us of the fact that the CMB constraint on \omh has a lensing component. The physics is the same as the well-known effect of shifting the reionization optical depth $\tau$: in the CMB, $A_s e^{-2\tau}$ controls the lensing amplitude, so a high-$\tau$ universe and a high-$A_\mathrm{lens}$ universe are nearly degenerate \citep{Sailer2026}. In their analysis \citet{Sailer2026} make this explicit: dropping Planck's low-$\ell$ polarization and letting the rest of the CMB$+$BAO data float $\tau$ within \LCDM yields $\tau = 0.090 \pm 0.012$, $\sim 3$--$5\sigma$ above the Planck low-$\ell$ value. Since essentially all current cosmological analyses simply adopt the Planck $\tau$ measurement, a systematic in that single number would propagate identically into essentially every CMB-anchored cosmological analysis.

\textbf{Relation to previous work.} Several recent analyses have examined the DESI results from angles closely related to ours. \citet{Efstathiou2025_BAO} rotates the BAO distance measurements into directions parallel and perpendicular to the Planck \LCDM degeneracy. The parallel direction effectively measures $\omega_m$, showing a ${\sim}2.1\sigma$ tension with Planck---in quantitative agreement with our $+2.2\sigma$ BAO \czero tension. The perpendicular direction and the equation of state at the pivot redshift, $w(z_{\rm piv} = 0.5) = -0.996 \pm 0.046$, show no departure from \LCDM, consistent with our finding that the dark energy direction \cone is null ($+1.2\sigma$). \citet{Efstathiou2025_SN} has also identified potential photometric systematics in the DES~Y5 supernova compilation, finding a ${\sim}0.04$~mag offset between low and high redshift subsamples; we use the recalibrated DES-Dovekie sample \citep{Popovic2025} which partially addresses these concerns. Our DES-Dovekie \czero tension is indeed the mildest of the three SN datasets ($-0.8\sigma$) but none is significant. 

 
\citet{Ye2025} independently confirm the ${\sim}2\sigma$ BAO--CMB tension using Bayesian suspiciousness, and find that the phantom crossing redshift $z_{\rm pc} = 0.45 \pm 0.04$ coincides with the best-constrained scale---matching our BAO \cone pivot $z = 0.46$. The pivot redshift approach was pioneered by \citet{CortesLiddle2024} in the context of DESI~DR1, who identified what they term the ``PhantomX coincidence'': that $w$ crosses $-1$ precisely at the epoch of observation. Our SVD analysis provides an explanation for this coincidence: the apparent $w \neq -1$ at the \czero pivot is \omh tension dressed in dark energy clothing, while the genuinely independent dark energy direction (\cone) is consistent with a cosmological constant at its own pivot. Note that \czero, although where the tension lives, is a near-perfect proxy for \omh --- a parameter that the CMB also constrains tightly along the same combination --- so the tension can in principle be resolved by any mechanism that shifts the CMB inference of \omh (e.g., $\tau$, $A_\mathrm{lens}$; Section~\ref{sec:extensions}). The dark-energy direction \cone, by construction, cannot be moved by such mechanisms --- which makes its null tension at $+1.2\sigma$ a genuine constraint on dynamical dark energy.

Our $V_0$ direction is the empirical, multi-probe counterpart of the analytically constructed Model-independent Expansion Discrepancy Index (MEDI) of \citet{Weiner2026_MEDI}: both single out the same parameter combination from very different starting points. The MEDI identifies $\omega_m r_d^2$ as the combination that acoustic-scale measurements constrain; within \LCDM, where $r_d$ is nearly fixed by early-universe physics, this reduces to \omh, which is why \czero is almost purely an \omh mode. The physical origin of the \omh dominance in BAO was recognized by \citet{Eisenstein2004} and \citet{Knox2006}. Appendix~\ref{app:universality} provides our own derivation connecting the log derivatives of distance observables to the $\beta$-vector coefficients. 

The curvature analysis in Section~\ref{sec:curvature} connects to \citet{ChenZaldarriaga2025}, who find $\Omk = +0.0023 \pm 0.0011$ from standard MCMC methods. We did not run a dedicated $\Omega_k$ MCMC chain in this work; we use simple formulas to estimate the value of $\Omega_k$ that the BAO data prefer. 

Finally, \citet{Sailer2026} argue that a systematic in Planck's low-$\ell$ polarization $\tau$ would act as a single-point failure for essentially all current CMB analyses, and that the rest of the CMB$+$BAO data within \LCDM prefer $\tau = 0.090 \pm 0.012$, a $\sim 3$--$5\sigma$ shift from the low-$\ell$ measurement. Their $\tau$ channel and our $A_\mathrm{lens}$ channel involve the same physics expressed differently: $A_s e^{-2\tau}$ controls the CMB lensing amplitude, so $A_\mathrm{lens}$ and $\tau$ are partially degenerate. 

The paper is organized as follows. Section~\ref{sec:data_method} describes the data and SVD methodology. Section~\ref{sec:universal_c0} establishes the universality of \czero across all probes. Section~\ref{sec:c0_tensions} presents the \czero tension measurements. Section~\ref{sec:w0wa_reinterpretation} reinterprets the \wowa preference. Section~\ref{sec:curvature} examines spatial curvature. Section~\ref{sec:extensions} considers other CMB-side extensions. Section~\ref{sec:conclusions} discusses the implications and summarizes our conclusions. Appendix~\ref{app:universality} provides an analytical derivation of the \czero universality.

This paper, both calculations and draft, was written using Claude Code. It is based on some old calculations done at the time that the original DESI paper appeared and partially presented in a talk at Perimeter Institute\footnote{https://pirsa.org/25080017}. The goal of those calculations was just for me to better understand the claims and are presented here in case anyone finds them interesting. Creating this manuscript served as a guide for me to understand how to use the new AI tools. I make no claims of originality or novelty.  All code, data-loading scripts, and the full analysis pipeline needed to reproduce every number and figure in this paper are available in a public GitHub repository.\footnote{\url{https://github.com/matiaszaldarriaga/desi-w0wa-svd}} It contains the SVD analysis package, the scripts that turn the public DESI~DR2 BAO, supernova, and CMB MCMC chains into the distance-ratio decompositions used here, the marimo notebook that regenerates every figure, and a machine-readable registry of all claims and quoted numbers; the aim is that a reader can re-derive the results from the public inputs, inspect each numerical choice, and freely modify or extend the analysis. This manuscript will not be submitted to any journal. It has been sufficiently ``refereed" by Claude Code\footnote{This is a joke, but I am definitely of the camp that wants to let LLMs ``cook" (https://arxiv.org/abs/2602.10181)}. I am happy to correct any mistakes and of course anyone is free to download the code and ask any agent to fix or improve it. Hopefully this disclaimer is sufficient to appease the arxiv police\footnote{https://xcancel.com/tdietterich/status/2055055541542760694}, but maybe transparency will trigger a ban and you will have to wait a year for any new paper with my name on it.